%%% LaTeX Template: Two column article
%%%
%%% Source: http://www.howtotex.com/
%%% Feel free to distribute this template, but please keep to referal to http://www.howtotex.com/ here.
%%% Date: February 2011

%%% Preamble
\documentclass[	DIV=calc,%
							paper=a4,%
							fontsize=12pt,%
							onecolumn]{scrartcl}	 					% KOMA-article class

\usepackage{lipsum}													% Package to create dummy text
\usepackage[english]{babel}										% English language/hyphenation
\usepackage[protrusion=true,expansion=true]{microtype}				% Better typography
\usepackage{amsmath,amsfonts,amsthm}					% Math packages
\usepackage[pdftex]{graphicx}									% Enable pdflatex
\usepackage[svgnames]{xcolor}									% Enabling colors by their 'svgnames'
\usepackage[hang, small,labelfont=bf,up,textfont=it,up]{caption}	% Custom captions under/above floats
\usepackage{epstopdf}												% Converts .eps to .pdf
\usepackage{subfig}													% Subfigures
\usepackage{booktabs}												% Nicer tables
\usepackage{fix-cm}													% Custom fontsizes
\usepackage[utf8]{inputenc}
\usepackage[top=2.5cm, bottom=2.5cm, left=2.5cm, right=2.5cm]{geometry}
\usepackage[ddmmyyyy]{datetime}
\addto\captionsenglish{%
	\renewcommand\tablename{Tabela}
	\renewcommand\figurename{Figura}
} 
 

 
%%% Custom sectioning (sectsty package)
\usepackage{sectsty}													% Custom sectioning (see below)
\allsectionsfont{%															% Change font of al section commands
	\usefont{OT1}{phv}{b}{n}%										% bch-b-n: CharterBT-Bold font
	}

\sectionfont{%																% Change font of \section command
	\usefont{OT1}{phv}{b}{n}%										% bch-b-n: CharterBT-Bold font
	}



%%% Headers and footers
\usepackage{fancyhdr}												% Needed to define custom headers/footers
	\pagestyle{fancy}														% Enabling the custom headers/footers
\usepackage{lastpage}	

% Header (empty)
\lhead{}
\chead{}
\rhead{}
% Footer (you may change this to your own needs)

%% ====================================
%% ====================================
%% mude o rodape  do projeto
%% ====================================
%% ====================================

\lfoot{\footnotesize \texttt{Projeto de Redes} \textbullet Reestruturação do Cabeamento}


\cfoot{}
\rfoot{\footnotesize página \thepage\ de \pageref{LastPage}}	% "Page 1 of 2"
\renewcommand{\headrulewidth}{0.0pt}
\renewcommand{\footrulewidth}{0.4pt}



%%% Creating an initial of the very first character of the content
\usepackage{lettrine}
\newcommand{\initial}[1]{%
     \lettrine[lines=3,lhang=0.3,nindent=0em]{
     				\color{DarkGoldenrod}
     				{\textsf{#1}}}{}}



%%% Title, author and date metadata
\usepackage{titling}															% For custom titles

\newcommand{\HorRule}{\color{DarkGoldenrod}%			% Creating a horizontal rule
									  	\rule{\linewidth}{1pt}%
										}

\pretitle{\vspace{-30pt} \begin{flushleft} \HorRule 
				\fontsize{50}{50} \usefont{OT1}{phv}{b}{n} \color{DarkRed} \selectfont 
				}

%% ====================================
%% ====================================
%% mude o titulo  do projeto
%% ====================================
%% ====================================

\title{Reestruturação de Infraestrutura de Redes}					% Title of your article goes here

%% ====================================



\posttitle{\par\end{flushleft}\vskip 0.5em}

\preauthor{\begin{flushleft}
					\large \lineskip 0.5em \usefont{OT1}{phv}{b}{sl} \color{DarkRed}}
\author{Leandro Tavares, Sidnei Gobetti de Lima }  	% Author name goes here


\postauthor{\footnotesize \usefont{OT1}{phv}{m}{sl} \color{Black} 
					\\Universidade Tecnológica Federal do Paraná - Câmpus Cornélio Procópio 								% Institution of author
					\par\end{flushleft}\HorRule}

\date{}																				% No date




%%% Begin document
\begin{document}
\maketitle
\thispagestyle{fancy} 	
\thispagestyle{empty}		% Enabling the custom headers/footers for the first page 
% The first character should be within \initial{}




%% ====================================
%% ====================================
%% mude o resumo  do projeto
%% ====================================
%% ====================================
\initial{E}\textbf{ste projeto vem propôr uma reestruturação do cabeamento de redes já existente em uma organização fictícia. Buscaremos 
	definir e implantar as melhores práticas nessa nova estruturação, objetivando um melhor desempenho da rede, qualidade de comunicação e assegurar a confiabilidade na troca de informações na rede.}

%% ====================================
\begin{figure}
	\centering
	\includegraphics{utfpr}
\end{figure}

\vspace{3cm}
\centerline{\textit{\textbf{\today}}}

\clearpage
    \renewcommand*\listfigurename{Lista de figuras}
\listoffigures

\renewcommand*\listtablename{Lista de tabelas}
\listoftables




\clearpage
\renewcommand{\contentsname}{Sumário}
\tableofcontents
\clearpage

%% ====================================
%% ====================================
%% Inicio do texto
%% ====================================
%% ====================================
\section{Introdução}
Será realizado uma reestruturação da rede de uma organização fictícia. Essa empresa atua no ramo alimentício, atendendo as cidades de Cascavel e Curitiba, possuindo um total de 12 filiais distribuídas nessas regiões. Contudo, a reestruturação será realizada apenas em uma filial, situada em Cascavel. Essa filial possui 200 colaboradores atuando em diversas atividades.

O parque de informática dispõe de 50 computadores, 5 servidores físicos, 3 switches, 2 controladores para a rede wireless, bem como impressoras e outros periféricos. A organização possui seu datacenter lotado na cidade de Curitiba, dispondo de vários sistemas, incluindo o sistema de ERP, aonde é acessado por terminal service, por todas as filiais.

Os pontos de venda (PDV) realizam a comunicação diretamente com o servidor em Curitiba. O meio de interconexão utilizado para realizar toda a comunicação entre as filiais e o datacenter é fibra óptica, e a comunição interna por cabeamento de cobre.

A reestruturação proposta para a filial de Cascavel em questão, terá o intuito de readequação da estrutura de rede para certificação. Uma vez que o projeto inicial não atingiu os requisitos para a certificação da estrutura. Outro ponto que será incluído no projeto é a implantação de uma estrutura virtual redundante para servidor de arquivos, pois segue um modelo que não atende a necessidade atual da empresa. 
\subsection{Benefícios}
A reestruturação da infraestrutura de rede trará benefícios impactantes, tais como:
\begin{itemize}
	\item Desempenho da rede melhorado: A eliminação dos gargalos presentes na rede irá aumentar o desempenho da rede, agilizando ainda mais o trafego de dados.
	
	\item Confiabilidade na transferência das informações: Com maior conformidade da infraestrutura, as informações terão maior garantia de entrega ao destino, e assim, aumentando o QoS ( Quality of Service).
	
	\item Menor custo de manutenção: Estando com a rede certificada, o custo com manutenções no cabeamento serão praticamente extintos.
\end{itemize}
Em decorrência disso, colaboradores podem aumentar sua produtividade, pois os problemas existentes devido a interrupção da comunicação, relacionadas a infraestrutura de rede, serão extremamente reduzidas. Os clientes ficarão mais satisfeitos, e dessa forma, os serviços que serão entregues podem agregar maior valor ao negócio. 

Por conta de tudo isso, a organização terá benefícios de curto a longo prazo, principalmente, econômicos.


\section{Estado atual}
Aprente o estado atual da rede. Caso não tenha rede, desconsiderar esta seção.

Caso tenha rede, deixe claro:
\begin{itemize}
	\item os passivos de rede atuais:path panels, cabos, etc..;
	\item as principais reclamações dos usuários. Qual o principal motivo da reestruturação? Efetue uma pesquisa junto aos colaboradores para determinar quais problemas a rede apresenta.
	\item Observações. Analise a rede e verifique se há estruturas que não se enquadram nas normas ou que indicam suspeita de problemas.
\end{itemize}


\section{Usuários e Aplicativos}
Explique nesta seção os usuários atuais e o perfil de crescimento, se por exemplo, há estimativa na evolução da empresa no que tange a quantidade de usuários, pontos de redes, equipamentos.
 

\subsection{Usuários}
Crie uma relação da quantidade, perfil de usuários de seu projeto.

\subsection{Aplicativos}
Crie uma relação dos aplicativos e seus níveis críticos de uso.


\section{Estrutura predial existente}

Explique aqui a planta física dos prédios
Pode ser anexada, em escala ou não.

Deve conter uma descrição geral, indicando a possível distância entre os pontos de rede e restrições de instalação.

\section{Planta Lógica - Elementos estruturados}

\subsection{Estado atual}
Deve ter a planta atual, se for o caso

\subsection{Topologia}
Proposta futura, proposta após implantação.
Deve conter o diagrama da rede. Atente-se a redundância  e ligações truncadas.
Deve explicar todos termos e componentes utilizados nestas plantas. Por exemplo: entrance facility, work area, horizontal cabling, etc..

Todos os elementos das figuras devem ser explicados. 
Crie esboço da configuração dos racks e brackets. Explique cada um dos componentes. Você pode criar uma tabela contendo figuras dentro, ou criar uma tabela e incluí-la como imagem. Por exemplo, verifique a tabela \ref{tab1}.

\begin{table}[h!]
\centering
\caption{Exemplo de tabela explicativa}
\label{tab1}
\begin{tabular}{|l|l|l|}
\hline
\multicolumn{3}{|l|}{Figura na Tabela} \\ \hline
1        & Rack          & \includegraphics[scale=0.2]{fig1}        \\ \hline
2        & Rack 2        & \includegraphics[scale=0.2]{fig1}        \\ \hline
\end{tabular}
\end{table}

\subsection{Encaminhamento}
Eletrodutos, calhas, e qualquer material em que os cabos serão alojados/alocados.

\subsection{Memorial descritivo}

Relacione todos os equipamentos passivos que serão utilizados, tipo, fabricante, quantidade.

\subsection{Identificação dos cabos}

\section{Implantação}
Estabeleça um cronograma de implantação:
Remoção de equipamentos existentes (destino para descarte), instalação dos condutores, instalação dos cabos, 
identificação dos cabos, montagem dos racks, certificação, etc... Crie atividades e estabeleça o tempo de execução. Se for um projeto real, indique também quais os responsáveis pela execução do projeto e de cada uma das etapas.

Defina marcas (e padrões) e fornecedores se for o caso. Atenção a contratados e subcontratados para a realização das atividades. Estabeleça a responsabilidade de execução da atividade e também da validação dela.

Utilize algum software para gerear o cronograma. Excel,etc. O fundamental é dividir em etapas, descrever e estimar o tempo de cada uma delas.

Segue uma relação de ferramentas:
http://asana.com/, 
https://trello.com/, 
http://www.ganttproject.biz/, 
http://www.orangescrum.org/. 

\section{Plano de certificação}
Quais seriam as etapas para a certificação? 
Quais os locais e horários para execução da certificação na rede? Toda rede será certificada?
Como os testes seriam executados?
Quais relatórios de certificação serão (ou deveriam ser) entregues? 

\section{Plano de manutenção}

Revisões periódicas na rede, emissão de certificados para novos pontos.

\subsection{Plano de expansão}
Existe um plano de expansão? Quantos novos pontos poderão ser acrecidos na rede, antes de migração de equipamentos na camada 2? Se houver expansão, quais equipamentos deverão ser direcionados para as estremidades da rede? 


\section{Orçamento}
Crie uma relação de orçamentos baseado na seções anteriores.

\section{Referências bibliográficas}
Utilize o mendley, o jabref ou diretamente o bibtex para gerenciar suas referências biliográficas. As referências são criadas automaticamente de acordo com o uso no texto.

Exemplo: Redes de computadores, segundo \cite{t2013} é considerada..... Já \cite{kurose2010} apresenta uma versão...

Analisando os pressupostos de \cite{ref3} e \cite{ref4} concluimos que....


\renewcommand\refname{} %%Referências bibliográficas}  
\bibliographystyle{ieeetr}
\bibliography{referencias}  

%% ***********************************************************************
%% === remover daqui =====================================================
%% ***********************************************************************
..
%% ***********************************************************************
%% === ate aqui    =====  ================================================
%% ***********************************************************************
\end{document}